\documentclass[11pt,a4paper]{article}
\newcommand{\docshort}{How to read the software}
\usepackage[utf8]{inputenc}
\usepackage[T1]{fontenc}
\usepackage{lmodern}
\usepackage{microtype}
\usepackage[margin=1in]{geometry}
\usepackage{booktabs}
\usepackage{tabularx}
\usepackage{enumitem}
\usepackage{xcolor}
\usepackage{fancyhdr}
\usepackage{listings}
\usepackage{hyperref}

\definecolor{codebg}{RGB}{245,247,250}
\definecolor{framecolor}{RGB}{210,220,230}
\definecolor{primary}{RGB}{20,50,90}

\hypersetup{colorlinks=true,linkcolor=primary,urlcolor=primary}

\pagestyle{fancy}
\fancyhf{}
\rhead{\small\textcolor{gray}{\docshort}}
\cfoot{\thepage}
\renewcommand{\headrulewidth}{0.4pt}

\newcommand{\reporoot}{..}
\newcommand{\code}[1]{\texttt{\detokenize{#1}}}
\newcommand{\coderef}[2]{\texttt{\detokenize{#1}}:{#2}}
\sloppy

\lstdefinestyle{src}{
  basicstyle=\footnotesize\ttfamily,
  backgroundcolor=\color{codebg},
  frame=single,
  rulecolor=\color{framecolor},
  numbers=left,
  numberstyle=\tiny\color{gray},
  numbersep=6pt,
  breaklines=true,
  breakatwhitespace=true,
  showstringspaces=false,
  tabsize=4,
  keepspaces=true,
  language=Python,
  keywordstyle=\color{primary}\bfseries,
  commentstyle=\color{gray}\itshape,
  stringstyle=\ttfamily
}
\lstset{style=src}
\setlist[itemize]{noitemsep,topsep=0.3em}
\setlist[enumerate]{noitemsep,topsep=0.3em}

\hypersetup{pdftitle={How to read the software}}

\title{\textbf{How to read the software}}
\author{fe-large-displacement}
\date{}

\begin{document}
\maketitle
\thispagestyle{fancy}

\noindent
Abaqus is several programs that share a name. This repository is a thin layer
on top of them. Confusing the parts is how a three-day job is spent on a
keyword CAE never wrote.

\section{Abaqus/CAE versus the solvers}

Abaqus/CAE builds a model. It does not solve it. Two solvers do:

\begin{itemize}
  \item Abaqus/Standard --- implicit. Equilibrium iterations each increment.
  \item Abaqus/Explicit --- explicit dynamics. A stable time increment, no
        iteration.
\end{itemize}

CAE writes an input file (\texttt{.inp}). The solver reads that file. CAE
accepts arguments it does not implement and omits keywords whose values match
the default. The deck is the ground truth. Every model here writes the
\texttt{.inp}, checks it, and only then submits.

\section{Standard versus Explicit}

\begin{center}
\small
\begin{tabular}{@{}lll@{}}
\toprule
 & Standard & Explicit \\
\midrule
Procedure & \texttt{*Static} (here) & \texttt{*Dynamic, Explicit} \\
Failure mode & no convergence & element distortion \\
User material & UMAT & VUMAT \\
Pore-pressure elements & \texttt{C3D8P} & none on ordinary solids \\
ALE & limited & production tool \\
CEL & none & Explicit only \\
Precision & double & set \texttt{DOUBLE\_PLUS\_PACK} \\
\bottomrule
\end{tabular}
\end{center}

Standard is the right tool until the soil starts to fail in a large-displacement
sense. Explicit is the right tool after that, provided the run is shown to be
quasi-static (\texttt{ALLKE}/\texttt{ALLIE}).

UMAT and VUMAT are different calling conventions for the same constitutive
idea. A closed-form Mohr--Coulomb return can be wrapped. A model with
substepping cannot be wrapped blindly. Verify on one element
(\texttt{constitutive/single\_element}) before a penetration job.

\section{ALE and CEL are solver features}

They are not separate products.

ALE keeps the mesh topology and moves nodes relative to the material. Explicit
ALE is what this repository uses in production. Standard ALE exists on
geometrically nonlinear \texttt{*Static} steps. It has one smoothing algorithm,
no initial mesh sweeps, and Lagrangian domains only. The implicit-ALE model
shows that path. It is not the tool for skirt penetration.

CEL lets material flow through a fixed Eulerian mesh. The structure stays
Lagrangian. CEL is Explicit only. There is no implicit CEL.

\section{What this repository is}

\begin{center}
\small
\begin{tabular}{@{}lp{0.62\textwidth}@{}}
\toprule
\texttt{examples/} & one complete script per formulation \\
\texttt{constitutive/} & UMAT/VUMAT sources, each with its own licence \\
\texttt{teaching/} & this PDF, and the formulations companion \\
\bottomrule
\end{tabular}
\end{center}

\section{How a job is launched}

\begin{verbatim}
set FLD_WORKDIR=C:\abq\run
set FLD_SUBMIT=0
abaqus cae noGUI=examples/implicit.py
\end{verbatim}

\texttt{FLD\_SUBMIT=0} writes the deck without solving. Start there.

When a user subroutine is compiled, launch through \texttt{abaqus.bat}, never
\texttt{ABQLauncher.exe} directly. Paths with spaces break the Fortran compile.

\section{The habit that catches silent failures}

Write the deck. Read it. Then submit. Three failures produce a normal ODB and
no warning: an ALE domain that never sweeps, a CEL domain with no void for
heave, and a VUMAT returning NaN (Explicit integrates it to the end).

On a new machine, write the decks with \texttt{FLD\_SUBMIT=0} before
submitting. The code that changes from implicit to CEL is
\texttt{teaching/formulations.pdf}.

\section{Clausen Mohr--Coulomb}

Use of \texttt{constitutive/mohr-coulomb-clausen/} requires citation of all
three papers. BibTeX is in that directory as \texttt{references.bib}.

\begin{enumerate}
  \item Clausen, J., Damkilde, L. and Andersen, L. (2006). Efficient return
        algorithms for associated plasticity with multiple yield planes.
        \emph{International Journal for Numerical Methods in Engineering}
        66(6), 1036--1059. \url{https://doi.org/10.1002/nme.1595}
  \item Clausen, J., Damkilde, L. and Andersen, L. (2007). An efficient return
        algorithm for non-associated plasticity with linear yield criteria in
        principal stress space. \emph{Computers \& Structures} 85(23),
        1795--1807. \url{https://doi.org/10.1016/j.compstruc.2007.04.002}
  \item Clausen, J., Damkilde, L. and Andersen, L.~V. (2015). Robust and
        efficient handling of yield surface discontinuities in elasto-plastic
        finite element calculations. \emph{Engineering Computations} 32(6),
        1722--1752. \url{https://doi.org/10.1108/EC-01-2014-0008}
\end{enumerate}

\end{document}
